\section{Definability}
Here we render explicit the standard (Kripke) version of our definability theorem, explain our choice of semantics, 
and suggest how we will transport the standard result to the new setting.

\subsection{Polyhedral semantics}
As is always the case, we begin with a syntax: here we consider the standard propositional modal language.
\begin{definition}[Modal language]
    Our language is the standard modal language: 
    we assume a set of propositional variables $X$ 
    and define inductively
    \[\Ll = \set{x\in X \ \|\ \bot\ \|\ \varphi\vee\psi\ \|\ \lnot\varphi\ \|\ \d\varphi}\]
\end{definition}
This is a minimal description of our syntax, 
where we can assume the ``not'' to be primitive and derive 
\begin{itemize}
    \item $\top=\lnot\bot$
    \item $\varphi \wedge \psi=\lnot(\lnot \varphi\vee\lnot\psi)$
    \item $\varphi\to\psi=\lnot \varphi\vee\psi$
    \item $\b \varphi=\lnot\d\lnot\varphi$
\end{itemize}
Our semantics is a tamer version of the topological 
semantics for modal logic (\cite{Gabelaia2001ModalDI}), 
interpreting the box $\b$ as an interior operator
and the diamond $\d$ as a closure operator.

The main difference between polyhedral and topological semantics 
is that the evaluations of the elementary propositions are 
(possibly open) subpolyhedra and not arbitrary subsets.

Recall that a subpolyhedron (\ref{Polyhedron}) is a union of relative interiors of subsimplices.

\begin{definition}[Polyhedral semantics]
    Let $P$ be a (compact) polyhedron. An \textbf{evaluation} is a map $\nu :\Ll \to \Pp(P)$ such that $\nu(\varphi)$ is a union of relative interiors of subpolyhedra of $P$.

    Moreover, we ask that $\nu$ satisfies the following properties:
    \begin{itemize}
        \item $\nu(\bot)=\emptyset$,
        \item $\nu(\varphi\vee\psi)=\nu(\varphi)\cup\nu(\psi)$,
        \item $\nu(\lnot\varphi)=P\setminus \nu(\varphi)$,
        \item $\nu(\d\varphi)=\overline{\nu(\varphi)}$.
    \end{itemize}
    We say $P,x\models \varphi$ if and only if $x\in \nu(\varphi)$, 
    and we say that $\varphi$ \textbf{is valid in} $P$, 
    written $P\models\varphi$, if $P,x\models \varphi$ 
    for every $x\in P$ and every evaluation $\nu$.
\end{definition}

This allows us to define the \textbf{logic} of a polyhedron, i.e. the set of formulas that are valid in a given polyhedron.
\begin{definition}
    Let $P$ be a polyhedron. We call the \textbf{logic} of $P$ the set
    \[\Log(P):=\set{\varphi\in \Ll:P\models\varphi}\]
\end{definition}

Recall that the Kripke semantics for the modal logic (S4.Grz) interprets the Boolean connectives as usual, together with
\[P,p\models \d\varphi\iff \exists q\geq p: P,q\models\varphi.\]
($\d \varphi$ is true at $p$ if and only if $\varphi$ is true on some element above $p$).

The following theorem connects the two semantics:
\begin{theorem}
    Let $P$ be a polyhedron and let $K=(V,\Sigma)\in \Tr(P)$ be a simplicial complex realizing $P$.

    For each $x\in P$ there exists a unique $\sigma_x\in \Sigma$ such that $x\in \RelInt(\sigma_x)$, and
    \[P,x\models\varphi \text{ (as polyhedra) if and only if }\Sigma,\sigma_x\models\varphi\text{ (as posets)}.\]
    Moreover,
    \[\Log(P)=\bigcap_{K\in \Tr(P)}\Log([K]).\]
\end{theorem}
\begin{proof}
    See \cite{TENCATE} and \cite{DBLP} for a complete account.
\end{proof}

In the form of a slogan: a formula is satisfied on a polyhedron if and only if it's satisfied on some triangulation,
and a formula is valid on a polyhedron if and only if it's valid on every triangulation.

\subsection{Modally definable classes}

The dual question to validity is definability: given a set of formulas $F$, we can define the class of models validating it:

\begin{definition}[Modally definable class]
    \label{MD1}
    Let $F\subseteq \Ll$ be a set of formulas. We call $\Pol(F)$ the class of (compact) polyhedra modeling every formula of $F$, i.e.
    \[\Pol(F)=\set{P\in\Pol: \forall\varphi\in F. P\models\varphi}\]
\end{definition}

Goldblatt and Thomason's theorem gives necessary and sufficient condition for when an arbitrary class of frames $\C$ is modally definable, i.e. 
it is the class of models validating some set of formulas.

Every flavor of models generates a slightly different Goldblatt-Thomason-like theorem. The classical result about 
frames states that a class of Kripke frames is modally definable if and only if it's closed under disjoint unions,
generated subframes, p-morphic images and it's closed under and reflects ultrafilter extensions.

Here we recall
\begin{definition}[Generated subframe]
    Let $(W,R)$ be a frame, $V\subseteq W$ is called a \textbf{generated subframe}
    if and only if it's upward-closed, i.e. for any $v\in V, $ if $vRv'$ then $v'\in V$ too.

    When $R=\leq$ this translates to the usual notion of upward closure.
\end{definition}

As we are considering compact polyhedra, our frames (posets) are finite, thus we can use a stronger theorem with easier conditions.

\begin{theorem}[Goldblatt-Thomason for finite transitive frames (I)]
    \label{GTFT1}
    A class of finite transitive frames $\C$ is modally definable if and only if it is closed under
    \begin{enumerate}
        \item finite disjoint unions,
        \item generated subframes,
        \item p-morphic images.
    \end{enumerate}
\end{theorem}
\begin{proof}
    See \cite{Blackburn_Rijke_Venema_2001}, Chapter 3.
\end{proof}

We see that generated subframes correspond to open subpolyhedra. This constitute a problem: Our polyhedra lay in some euclidean space, thus have open non-compact subsets.

Such a subset cannot be part of our class as it's not compact, but a priori there might be maps that can be defined on it and not on the starting space.

Our solution is to avoid considering generated subframes (equivalently open subpolyhedra) and combine the notion of generated subframe and p-morphic image into one notion.

We need to show equivalence, so let us state it as a corollary and prove it.
\begin{corollary}[Goldblatt-Thomason for finite transitive frames (II)]
    A class of finite transitive frames $\C$ is modally definable if and only if it is closed under
    \begin{enumerate}
        \item finite disjoint unions,
        \item[(3')] surjective up-reductions,
    \end{enumerate}
\end{corollary}
\begin{proof}
    Suppose $\C$ is closed under the conditions (1),(2), and (3) of \ref{GTFT1}. Let $X\in \C$ and let $f:X\upred Y$ be an up-reduction. As an up-reduction is a partial function, call $U$ 
    the domain of definition of $f$. Thanks to condition (2) we know that $U\in \C$.

    Then $f_{|U}: U\twoheadrightarrow Y$ is a surjective p-morphism, 
    so $Y$ is a p-morphic image of $U$, and thus by (3) $Y\in \C$. 
    This means that $\C$ is closed under (3').

    Conversely, suppose that $\C$ is closed under (1) and (3'). 
    Let $X\in \C$ and let $U$ be a generated subframe of 
    $X$. Then the identity on $U$, 
    seen as a partial map $X\to U$, is an up-reduction. 
    Similarly, if $Y$ is a p-morphic image of $Y$, then, 
    since $X$ is a generated subframe of $X$, we may extract an up-reduction $X\upred Y$.
\end{proof}

Note that this means the class of frames is always closed under generated subframes, but we may consider this class within a category that lacks them.

We will also make use of the Jankov-Fine axiomatization: There exists some formulas that prohibit up-reduction to a given (rooted) frame, i.e. a model validates 
the formula corresponding to the frame if and only if it does not up-reduce to the frame.

These formulas are useful for example and counterexamples and fully characterize classes of finite frames.

\begin{theorem}[Jankov-Fine]
    Let $U$ be a finite rooted frame. Then there exists a formula $\chi(U)$ such that $F\models\chi(U)$ if $F$ does \textbf{not} up-reduce to $U$.
\end{theorem}
\begin{proof}
    See \cite{Bezhanishvili2022}.%provide a more direct source?
\end{proof}
\begin{corollary}
    A class of finite frames $\C$ can be characterized as
    \[\C=\Mod(\set{\chi(U)\ |\ U\notin \C,\ U \text{ rooted}}).\]
\end{corollary}

